\section{Plugin CTFd}\label{plugin-ctfd}

CTFd è una delle piattaforme open source più diffuse per la gestione di
competizioni Capture The Flag. Nativamente, CTFd è progettato per sfide
di tipo \emph{Jeopardy}, in cui la risoluzione coincide con l'invio di
una stringa alfanumerica statica o regex (la \emph{flag}).

Per abilitare gare di tipo \emph{Defense} o \emph{Patching}, in cui
l'obiettivo del team consiste nel mettere in sicurezza un servizio
vulnerabile mantenendone la piena operatività, è stata sviluppata
un'estensione per CTFd composta da due moduli integrati:
\passthrough{\lstinline!patch\_challenges!} e
\passthrough{\lstinline!patch\_flags!}. Il plugin agisce da
intermediario tra l'interfaccia web della competizione e le API REST
esposte da BlueAgent.

\begin{quote}
CTFd usa Flask al posto di FastAPI (adottato in BlueAgent) ed usa
SQLAlchemy per gestire facilmente il modello dei dati utilizzato nel
codice.
\end{quote}

\subsection{Componenti del Plugin: PatchChallenge e
PatchFlag}\label{componenti-del-plugin-patchchallenge-e-patchflag}

L'integrazione sfrutta l'architettura a plugin di CTFd estendendone le
classi base:

\begin{itemize}
\tightlist
\item
  \textbf{\passthrough{\lstinline!PatchChallenge!}}: definisce una nuova
  tipologia di sfida. Ha un modello dei dati per facilitare
  l'integrazione con il database di CTFd. Questo modello memorizza la
  specifica dichiarativa della topologia in formato JSON/YAML e
  l'insieme dei comandi di verifica da inoltrare a BlueAgent.
\item
  \textbf{\passthrough{\lstinline!PatchFlag!}
  (\passthrough{\lstinline!BaseFlag!})}: invece di confrontare una
  stringa inserita dal giocatore con un valore memorizzato nel database,
  il metodo \passthrough{\lstinline!compare()!} interroga l'Oracolo
  (BlueAgent) e registra il tentativo di verifica ed il suo esito.
\end{itemize}

\subsection{Distribuzione delle Configurazioni
WireGuard}\label{distribuzione-delle-configurazioni-wireguard}

Per permettere ai partecipanti di operare sui container della propria
challenge, il plugin gestisce automaticamente l'accesso alla rete del
team:

\begin{enumerate}
\def\labelenumi{\arabic{enumi}.}
\tightlist
\item
  Il plugin espone all'interfaccia web le rotte per la gestione della
  connettività.
\item
  Quando un utente autenticato visualizza la sfida o richiede l'accesso,
  CTFd invoca
  \passthrough{\lstinline!PUT /teams/\{team\_k\}/members/\{user\_id\}!}
  di BlueAgent.
\item
  BlueAgent alloca dinamicamente un indirizzo IP libero all'interno
  della subnet VPN del team (\passthrough{\lstinline!10.k.0.0/24!}),
  genera la coppia di chiavi asimmetriche e registra il peer
  sull'interfaccia WireGuard.
\item
  Il plugin restituisce al browser il file di configurazione WireGuard
  (\passthrough{\lstinline!.conf!}) pronto per l'importazione.
\end{enumerate}

\subsection{Flusso di Verifica e
Convalida}\label{flusso-di-verifica-e-convalida}

Il ciclo di verifica e attribuzione del punteggio si articola nelle
seguenti fasi:

\begin{enumerate}
\def\labelenumi{\arabic{enumi}.}
\tightlist
\item
  \textbf{Richiesta di verifica}: il giocatore, dopo aver applicato le
  contromisure sui file e sui servizi del container, preme il pulsante
  di verifica sull'interfaccia di CTFd.
\item
  \textbf{Inoltro all'Oracolo}: CTFd invia una richiesta
  \passthrough{\lstinline!POST /teams/\{team\_k\}/verificator!} a
  BlueAgent, trasmettendo l'identificativo del team e la specifica della
  challenge.
\item
  \textbf{Esecuzione}: BlueAgent esegue la verifica in ambiente
  separato.
\item
  \textbf{Aggiornamento punteggio}: BlueAgent restituisce il verdetto a
  CTFd che si occupa di aggiornare i punteggi.
\end{enumerate}
